\chapter{Glosario operacional}

Este glosario fija el uso de términos dentro del manuscrito. No pretende reemplazar
definiciones generales de ingeniería de software.

\begin{longtable}{>{\bfseries}p{3.6cm}p{10.6cm}}
\caption{Glosario COBRIA.}\label{tab:glosario-cobria}\\
\toprule
Término & Definición operacional \\
\midrule
\endfirsthead
\toprule
Término & Definición operacional \\
\midrule
\endhead
Autoridad & Capacidad reconocida para establecer un hecho de negocio.\\
Canónico & Representación autoritativa de una entidad dentro de un ámbito y versión.\\
Catálogo & Vocabulario o conjunto de referencias permitidas y gobernadas.\\
Conector & Adaptador que encapsula la API de almacenamiento o proveedor.\\
Contrato & Especificación versionada de entradas, salidas, políticas y verificaciones.\\
Dater & Nombre local de una clase que combina acceso, mapeo, Repositorio, caché o índices.\\
Conjunto de datos & Proyección analítica versionada con selección, transformación y manifiesto.\\
Derivado & Artefacto calculado desde una fuente que no adquiere autoridad automáticamente.\\
Entidad & Objeto con identidad, ámbito, atributos, ciclo de vida, versión y revisión.\\
Registro de evidencias & Registro minimizado que enlaza solicitud, evidencia, decisión y efecto.\\
Conjunto de características & Conjunto versionado de características con tiempo y procedencia.\\
Idempotencia & Propiedad por la cual repetir una operación no multiplica su efecto.\\
Inferencia & Salida probabilística de un modelo, distinta de un hecho confirmado.\\
Linaje & Grafo que conecta artefactos con fuentes, transformaciones y versiones.\\
Mapeador & Componente que traduce entre dominio y representación física.\\
Metadato & Dato que declara estructura, relación, política, versión o significado.\\
Cola de salida & Registro persistente de una intención de publicación asociada a un cambio.\\
Proyección & Representación derivada optimizada para una lectura, proceso o análisis.\\
Gestor de proyecciones & Componente que mantiene, verifica, reconstruye y versiona proyecciones.\\
Reconstruibilidad & Capacidad de regenerar un derivado desde fuentes y contratos suficientes.\\
Repositorio & Interfaz de acceso a entidades expresada en lenguaje del dominio.\\
Revisión & Contador o token utilizado para detectar actualización concurrente obsoleta.\\
Ámbito & Frontera de pertenencia y autorización, como organización, organización o contexto.\\
Servicio & Componente que expresa coordinación o capacidad de dominio/aplicación.\\
Caso de uso & Frontera autorizada que coordina una intención y sus efectos.\\
Marca de corte & Punto temporal o cursor declarado para procesar cambios o cortes.\\
\bottomrule
\end{longtable}

\chapter{Plantillas de contratos}
\begingroup
\small

\section{Contrato de entidad}

\begin{verbatim}
entity:
  name: ExampleEntity
  version: 1
  identity: id
  ámbito: [organizaciónId]
  revision: revision
  lifecycle: [draft, active, archived]
  fields:
    - name: title
      type: string
      required: true
      sensitivity: internal
  invariants:
    - "title must not be empty"
  retention: policy-entity-default
  responsable: domain-team
\end{verbatim}

\section{Contrato de Mapeador}

\begin{verbatim}
mapper:
  entity: ExampleEntity
  version: 2
  readsPhysicalVersions: [1, 2]
  writesPhysicalVersion: 2
  fields:
    id: key
    organizaciónId: organización_id
    title: display_name
    revision: rev
  roundTrip:
    equality: semantic
    ignored: [storage_updated_at]
\end{verbatim}

\section{Contrato de proyección}

\begin{verbatim}
projection:
  id: example-by-status
  version: 3
  sources: [ExampleEntity@1]
  key: [organizaciónId, status, id]
  consistency: eventual
  freshnessSloSeconds: 30
  deterministic: true
  idempotent: true
  authoritative: false
  rebuild:
    mode: [full, incremental]
    batchSize: 500
    checkpoint: required
    cambio de versión: version-alias
  verify:
    coverage: 1.0
    orphanCount: 0
  responsable: query-team
\end{verbatim}

\section{Contrato de conjunto de datos}

\begin{verbatim}
conjunto de datos:
  id: example-training-set
  version: 4
  purpose: approved-task
  ámbito: organización
  sources: [ExampleEntity@1, ExampleEvent@2]
  cutoff: required
  transformationsCommit: abc123
  esquema: example-conjunto de datos-esquema@4
  quality:
    completenessMin: 0.98
    duplicateMax: 0
  privacy:
    allowedFields: [category, eventTime, value]
    retentionDays: 90
  lineage: required
  responsable: equipo-analitica
\end{verbatim}

\section{Contrato de característica}

\begin{verbatim}
feature:
  name: rolling_count_30d
  version: 2
  entity: ExampleEntity
  type: integer
  window: 30d
  availabilityTime: ingestionTime
  toleranciaFueraEnLinea: 0
  freshnessSeconds: 300
  sourceDataset: example-events@2
  responsable: ml-platform
\end{verbatim}

\section{Contrato de herramienta de agente}

\begin{verbatim}
tool:
  name: propose_entity_change
  version: 1
  capability: entity.change.propose
  inputSchema: proposal-input@1
  outputSchema: proposal-output@1
  ámbito: inherited-and-verified
  risk: medium
  confirmation: required-before-effect
  idempotencyKey: required
  reversible: true
  compensation: cancel-proposal
  audit: evidence-ledger
\end{verbatim}

\endgroup
\chapter{Listas de comprobación de ingeniería y gobierno}

\section{Lista de comprobación de entidad canónica}

\begin{longtable}{p{1cm}p{12.8cm}}
\caption{Verificación de una entidad canónica.}\label{tab:check-entidad}\\
\toprule
OK & Criterio \\
\midrule
\endfirsthead
\toprule
OK & Criterio \\
\midrule
\endhead
$\square$ & Posee identidad estable y definida.\\
$\square$ & El ámbito es obligatorio y forma parte de toda clave de acceso.\\
$\square$ & Distingue atributos, relaciones y derivados.\\
$\square$ & Declara versión de esquema y estrategia de lectura histórica.\\
$\square$ & Usa revisión o precondición para conflictos.\\
$\square$ & Sus invariantes se ejecutan dentro de casos de uso.\\
$\square$ & No conoce rutas ni SDK del proveedor.\\
$\square$ & Tiene responsable, clasificación y retención.\\
$\square$ & Su Mapeador pasa ida y vuelta semántico.\\
$\square$ & Existe una sola autoridad para cada hecho.\\
\bottomrule
\end{longtable}

\section{Lista de comprobación de proyección}

\begin{longtable}{p{1cm}p{12.8cm}}
\caption{Verificación de una proyección reconstruible.}\label{tab:check-proyeccion}\\
\toprule
OK & Criterio \\
\midrule
\endfirsthead
\toprule
OK & Criterio \\
\midrule
\endhead
$\square$ & Declara problema de lectura y consumidores.\\
$\square$ & Identifica fuentes y versiones suficientes.\\
$\square$ & No acepta escrituras autoritativas.\\
$\square$ & Define ámbito, clave, función y versión.\\
$\square$ & Declara consistencia y frescura.\\
$\square$ & Es idempotente o documenta por qué no puede serlo.\\
$\square$ & Posee rebuild total; incremental si está justificado.\\
$\square$ & Verifica cobertura, huérfanos y diferencias.\\
$\square$ & Tiene presupuesto, responsable, observabilidad y retiro.\\
$\square$ & Puede hacer cambio de versión o reversión sin mezclar versiones.\\
\bottomrule
\end{longtable}

\section{Lista de comprobación multiinquilino}

\begin{longtable}{p{1cm}p{12.8cm}}
\caption{Verificación de aislamiento por ámbito.}\label{tab:check-organización}\\
\toprule
OK & Criterio \\
\midrule
\endfirsthead
\toprule
OK & Criterio \\
\midrule
\endhead
$\square$ & API, caso de uso y Repositorio requieren ámbito.\\
$\square$ & Rutas e índices incluyen ámbito antes de identidad.\\
$\square$ & Claves de caché y de idempotencia incluyen ámbito.\\
$\square$ & Procesos de trabajo vuelven a autorizar mensajes.\\
$\square$ & Exports, backups y conjuntos de datos conservan separación.\\
$\square$ & Reconstrucciones requieren alcance explícito.\\
$\square$ & Telemetría no expone cargas útiles de otros organizaciones.\\
$\square$ & IDs iguales entre organizaciones forman parte de las pruebas.\\
$\square$ & Las pruebas negativas abarcan lectura y escritura.\\
$\square$ & Una fuga se trata como fallo crítico, no como tasa promedio.\\
\bottomrule
\end{longtable}

\section{Lista de comprobación de IA y agentes}

\begin{longtable}{p{1cm}p{12.8cm}}
\caption{Verificación de componentes inteligentes.}\label{tab:check-ia}\\
\toprule
OK & Criterio \\
\midrule
\endfirsthead
\toprule
OK & Criterio \\
\midrule
\endhead
$\square$ & Conjunto de datos y Conjunto de características poseen propósito, versión y linaje.\\
$\square$ & El corte temporal evita utilizar información futura.\\
$\square$ & Modelo y instrucción están versionados.\\
$\square$ & Recuperación filtra ámbito y permisos antes de generar.\\
$\square$ & La evidencia enlaza fragmento y fuente vigente.\\
$\square$ & El contenido recuperado no puede redefinir políticas.\\
$\square$ & Herramientas tienen capacidades mínimas y schemas.\\
$\square$ & Acciones sensibles solicitan confirmación.\\
$\square$ & Las operaciones mutantes son idempotentes.\\
$\square$ & Una inferencia no se persiste como hecho sin caso de uso.\\
$\square$ & Existen monitoreo, retiro y respuesta a incidentes.\\
\bottomrule
\end{longtable}

\chapter{Diccionario de métricas}

\begin{longtable}{>{\bfseries}p{3.2cm}p{5.2cm}p{5.2cm}}
\caption{Definición de métricas del estudio.}\label{tab:diccionario-metricas}\\
\toprule
Métrica & Cálculo o unidad & Precaución \\
\midrule
\endfirsthead
\toprule
Métrica & Cálculo o unidad & Precaución \\
\midrule
\endhead
Latencia p50 & Mediana en ms por corrida & No sustituir por media.\\
Latencia p95/p99 & Percentil por corrida & Requiere muestra suficiente.\\
Rendimiento efectivo & Operaciones exitosas por segundo & Reportar errores y concurrencia.\\
$A_r$ & Bytes leídos / bytes útiles & Definir respuesta vacía.\\
$A_w$ & Escrituras físicas / escritura lógica & Separar réplicas internas.\\
$A_s$ & Almacenamiento total / canónico & Incluir índices y staging.\\
Frescura & Confirmación a visibilidad correcta & Usar reloj y frontera declarados.\\
Cobertura & Esperados presentes / esperados & No detecta registros extra.\\
Huérfanos & Derivados sin fuente válida & Reportar conteo absoluto.\\
Exactitud de reconstrucción & $1-(faltantes+extras+diferentes)/esperados$ & Desglosar errores.\\
Cobertura de linaje & Artefactos con ruta completa / evaluados & No implica calidad.\\
Deriva & Distancia o prueba por variable & Cambio no siempre es defecto.\\
Precisión@k & Relevantes recuperados / $k$ & Depende del juicio de relevancia.\\
Exhaustividad@k & Relevantes recuperados / relevantes & Requiere conjunto completo.\\
Abstención & Sin respuesta cuando no hay evidencia & Evaluar falsos rechazos.\\
Violaciones ámbito & Accesos o efectos cruzados & Objetivo obligatorio: cero.\\
Esfuerzo & Tiempo activo y artefactos modificados & Separar aprendizaje.\\
Tamaño de efecto & Diferencia estandarizada o razón & Informar intervalo.\\
\bottomrule
\end{longtable}

\chapter{Esquema del paquete reproducible}

\section{Estructura de directorios}

\begin{verbatim}
cobria-replication/
  README.md
  LICENSE
  CITATION.cff
  environment/
    versions.bloqueo
    hardware-template.json
  reference/
    src/
    tests/
  contracts/
    entities/
    projections/
    conjuntos de datos/
    tools/
  generators/
    profile-r/
    profile-a/
    profile-b/
  experiments/
    performance/
    reconstruction/
    isolation/
    evolution/
    intelligence/
  analysis/
    notebooks/
    scripts/
  results/
    crudo/
    processed/
    figures/
    manifests/
\end{verbatim}

\section{Manifiesto de corrida}

\begin{verbatim}
{
  "runId": "uuid",
  "startedAt": "ISO-8601",
  "revisión": "git-sha",
  "profile": "R|A|B",
  "configuration": "B0|B1|C1|C2|C3|C4",
  "seed": 42001,
  "N": 100000,
  "organizaciones": 10,
  "workload": "W2",
  "concurrency": 32,
  "projectionCount": 3,
  "environmentHash": "sha256",
  "conjunto de datosHash": "sha256",
  "result": "completed|failed|excluded",
  "exclusionReason": null
}
\end{verbatim}

\section{Formato de evento de medición}

\begin{verbatim}
{
  "runId": "uuid",
  "operationId": "uuid",
  "operation": "query-by-status",
  "organizaciónAlias": "t-03",
  "startedNs": 0,
  "durationNs": 0,
  "reads": 0,
  "writes": 0,
  "bytesRead": 0,
  "bytesWritten": 0,
  "retries": 0,
  "resultCount": 0,
  "status": "ok|business-reject|error|timeout"
}
\end{verbatim}

\section{Reglas de publicación}

Los resultados crudo son inmutables. El procesamiento escribe otra carpeta y conserva
versión del guion. Las tablas del manuscrito se generan, no se editan manualmente. Una
exclusión permanece en crudo y aparece en el reporte de flujo. Los secretos se suministran
fuera del paquete y nunca se imprimen.

\section{Lista de comprobación de réplica}

\begin{enumerate}[leftmargin=*]
  \item verificar hash y licencia del paquete;
  \item instalar versiones fijadas;
  \item registrar hardware y límites;
  \item generar perfiles con semillas publicadas;
  \item ejecutar pruebas de corrección antes del prueba comparativa;
  \item correr configuraciones en orden aleatorio;
  \item conservar registros y manifiestos;
  \item ejecutar análisis sin modificar crudo;
  \item comparar tablas y documentar divergencias;
  \item publicar entorno y resultados negativos.
\end{enumerate}

\chapter{Plantillas de registro y publicación}

\section{Ficha de experimento}

\begin{table}[H]
\centering
\caption{Plantilla de ficha experimental.}
\label{tab:ficha-experimento}
\begin{tabularx}{\textwidth}{p{3.8cm}X}
\toprule
Campo & Registro \\
\midrule
ID y versión & \rule{8cm}{0.4pt}\\
Hipótesis & \rule{8cm}{0.4pt}\\
Variable independiente & \rule{8cm}{0.4pt}\\
Variables dependientes & \rule{8cm}{0.4pt}\\
Controles & \rule{8cm}{0.4pt}\\
Conjunto de datos y hash & \rule{8cm}{0.4pt}\\
Configuraciones & \rule{8cm}{0.4pt}\\
Repeticiones & \rule{8cm}{0.4pt}\\
Exclusiones previas & \rule{8cm}{0.4pt}\\
Regla de decisión & \rule{8cm}{0.4pt}\\
Revisión & \rule{8cm}{0.4pt}\\
\bottomrule
\end{tabularx}
\end{table}

\section{Ficha de resultado}

\begin{table}[H]
\centering
\caption{Plantilla de resultado.}
\label{tab:ficha-resultado}
\begin{tabularx}{\textwidth}{p{3.8cm}X}
\toprule
Campo & Registro \\
\midrule
Corridas incluidas & \rule{8cm}{0.4pt}\\
Corridas excluidas & \rule{8cm}{0.4pt}\\
Resumen descriptivo & \rule{8cm}{0.4pt}\\
Tamaño de efecto & \rule{8cm}{0.4pt}\\
Intervalo & \rule{8cm}{0.4pt}\\
Resultado de hipótesis & Apoyada / no apoyada / inconclusa\\
Costo adverso & \rule{8cm}{0.4pt}\\
Amenaza principal & \rule{8cm}{0.4pt}\\
Artefactos & \rule{8cm}{0.4pt}\\
\bottomrule
\end{tabularx}
\end{table}

\section{Ficha de incidente experimental}

Un incidente registra identificador de corrida, instante, etapa, síntoma, impacto, datos afectados, causa
conocida, decisión de exclusión, evidencia y acción preventiva. Un incidente de
aislamiento se reporta aunque la corrida no sea válida para rendimiento.

\section{Declaración de disponibilidad}

La publicación deberá indicar qué artefactos son abiertos, cuáles se sustituyeron por
datos sintéticos y cuáles no pueden compartirse. La imposibilidad de publicar código
privado no impide compartir implementación neutral, generadores, contratos y resultados.

\section{Declaración de autoría}

\ifnum\blindreview=1
La identidad del investigador y autor se omite en esta copia para revisión ciega. Las contribuciones
\else
El investigador y autor de la propuesta es Billy Jak Cordero Porras. Las contribuciones
\fi
de revisores, colaboradores y herramientas se declararán según las políticas de la
institución y de la sede de publicación. La autoría implica responsabilidad sobre
afirmaciones, datos, análisis y correcciones.

\chapter{Matriz de trazabilidad del manuscrito}

\begin{longtable}{>{\bfseries}p{1.4cm}p{4.2cm}p{4.6cm}p{3.6cm}}
\caption{Trazabilidad entre preguntas, artefactos y evaluación.}\label{tab:trazabilidad-manuscrito}\\
\toprule
RQ & Artefacto COBRIA & Experimento & Resultado requerido \\
\midrule
\endfirsthead
\toprule
RQ & Artefacto COBRIA & Experimento & Resultado requerido \\
\midrule
\endhead
RQ1 & Modelo de complejidad & P1--P5, R1--R2 & Curvas y amplificación.\\
RQ2 & Proyecciones de consulta & P1 & Latencia, lecturas y bytes.\\
RQ3 & Presupuesto de proyección & P2--P3 & $A_w$, $A_s$ y equilibrio.\\
RQ4 & Gestor de proyecciones & R1--R5 & Exactitud, tiempo y recuperación.\\
RQ5 & Metadatos y generador & Evolución y generator test & Esfuerzo y divergencia.\\
RQ6 & Ámbito y políticas & Pruebas negativas & Cero fugas y efectos.\\
RQ7 & Perfiles A y B & Mapeo y ejecución comparada & Principios sin ruptura.\\
RQ8 & Conjunto de datos/Feature contracts & AI-1--AI-2 & Linaje y equivalencia.\\
RQ9 & Retrieval, tools y ledger & AI-3--AI-7 & Evidencia y acciones seguras.\\
RQ10 & Lienzo y costo & Todos, incluidos negativos & Región de no adopción.\\
\bottomrule
\end{longtable}

\section{Trazabilidad de partes}

\begin{table}[H]
\centering
\caption{Función de cada parte del documento maestro.}
\label{tab:trazabilidad-partes}
\begin{tabularx}{\textwidth}{p{2cm}p{4.3cm}X}
\toprule
Parte & Función & Producto \\
\midrule
I & Plantear problema y preguntas & Hipótesis y alcance.\\
II & Formalizar artefacto & Modelo, contratos y algoritmos.\\
III & Razonar sobre costo & Complejidad y optimización.\\
IV & Extender a datos e IA & Linaje, RAG, agentes y gobierno.\\
V & Conectar con implementaciones & Casos anónimos y referencia.\\
VI & Diseñar evaluación & Experimentos y análisis.\\
VII & Sintetizar y concluir & Guía, contribuciones y límites.\\
VIII & Facilitar réplica & Plantillas, checklists y diccionario.\\
\bottomrule
\end{tabularx}
\end{table}

\section{Criterio de completitud para una versión publicable}

El manuscrito estará listo para someterse como estudio empírico cuando: la referencia y
los generadores sean públicos; P1, P2, R1, R2, aislamiento, evolución, AI-1 y AI-3 tengan
corridas reproducibles; las hipótesis se actualicen sin eliminar resultados negativos;
otro evaluador revise una muestra; y tablas y figuras provengan de scripts.

Si se publica antes de esos experimentos, debe presentarse como propuesta arquitectónica
y protocolo de investigación, no como validación de rendimiento. Esta distinción debe
aparecer en título, resumen, conclusiones y carta de presentación.

\section{Cierre de los anexos}

Los anexos transforman COBRIA de una descripción conceptual en un artefacto inspeccionable.
Los contratos hacen explícitas sus propiedades; los checklists permiten auditar una
implementación; el diccionario fija métricas; y el paquete reproducible define cómo
convertir futuras ejecuciones en evidencia.

\enlargethispage{2\baselineskip}
La utilidad de estas plantillas depende de su aplicación crítica. Marcar todas las casillas
sin pruebas no aumenta madurez. Cada afirmación debe conservar una ruta hacia código,
configuración, datos, ejecución o revisión que otra persona pueda examinar.
